\item La función de onda del estado fundamental para el electrón en un átomo de hidrógeno es igual a:
$$ \psi_{1s}(r) = \frac{1}{\sqrt{\pi a_0^3}} e^{-r/a_0} $$
donde $r$ es la coordenada radial del electrón y $a_0$ es el radio de Bohr.
\begin{enumerate}
    \item Demuestre que la función de onda, como se ha dado, está normalizada.
    \item Determine la probabilidad de localizar al electrón entre $r_1 = a_0/2$ y $r_2 = 3a_0/2$.
\end{enumerate}